\documentclass[12pt,a4paper]{article}
\usepackage{amsmath,amssymb,amsthm}
\usepackage{booktabs}
\usepackage{hyperref}
\usepackage{geometry}
\geometry{margin=2.5cm}
\usepackage{enumitem}

\newtheorem{theorem}{Theorem}[section]
\newtheorem{lemma}[theorem]{Lemma}
\newtheorem{corollary}[theorem]{Corollary}
\newtheorem{proposition}[theorem]{Proposition}
\theoremstyle{definition}
\newtheorem{definition}[theorem]{Definition}
\theoremstyle{remark}
\newtheorem{remark}[theorem]{Remark}

\title{Complex Numbers from Recognition Science:\\
Why Physics Requires $\mathbb{C}$}

\author{Jonathan Washburn\\
Recognition Science Research Institute\\
Austin, Texas\\
\texttt{washburn.jonathan@gmail.com}}

\date{March 2026}

\begin{document}
\maketitle

\begin{abstract}
We derive the necessity of complex numbers in physics from the
8-tick epoch of Recognition Science (RS).  The ledger updates in
a cycle of $2^D = 8$ phases (for $D = 3$), generating the cyclic
group $C_8$.  A faithful representation of $C_8$ requires the
complex numbers: the real line $\mathbb{R}$ cannot distinguish
all eight phases.  The primitive 8th root of unity
$\omega = e^{i\pi/4}$ is the tick-advance operator, and the
imaginary unit $i = \omega^2$ is the quarter-period operator
satisfying $i^2 = -1$ because two quarter-periods produce the
half-period reversal.  Euler's identity $e^{i\pi} + 1 = 0$ is
the 8-tick cycle-completion identity: the half-period amplitude
$\omega^4 = -1$ plus the ground state $\omega^0 = 1$ equals
zero.  Complex numbers are not imported into physics as a
mathematical convenience---they are forced by the discrete
phase structure of the ledger.  All structural results are
machine-verified.
\end{abstract}

%% ============================================================
\section{Introduction}
\label{sec:intro}
%% ============================================================

Why does physics require complex numbers?  Quantum mechanics
uses complex amplitudes ($\psi \in \mathbb{C}$),
electromagnetism uses complex phasors, signal processing uses
the Fourier transform $e^{i\omega t}$, and the Dirac equation
lives in a complex spinor space.  Yet the real world is
observed through real-valued measurements.  The standard view
treats $\mathbb{C}$ as a mathematical convenience: complex
numbers simplify calculations, but one could in principle work
in $\mathbb{R}^2$.

Recognition Science offers a sharper answer: complex numbers
are \emph{forced} by the discrete phase structure of the
ledger.  The 8-tick epoch---the fundamental update cycle of
the RS ledger---generates a cyclic symmetry that cannot be
faithfully represented over $\mathbb{R}$.  The minimum field
extension that supports this symmetry is $\mathbb{C}$.

Recent work by Renou et al.~\cite{Renou2021} and
Procopio et al.~\cite{Procopio2022} has shown experimentally
that quantum mechanics cannot be formulated with real Hilbert
spaces alone, providing independent evidence for the necessity
of $\mathbb{C}$.

%% ============================================================
\section{The 8-Tick Epoch}
\label{sec:epoch}
%% ============================================================

\begin{definition}[8-Tick Phase Space]
\label{def:phases}
The RS ledger updates in cycles of
$2^D = 2^3 = 8$ ticks.  The phase space is the cyclic group
\begin{equation}
  \Phi = \mathbb{Z}/8\mathbb{Z} = \{0, 1, 2, 3, 4, 5, 6, 7\},
\end{equation}
with the tick-advance operation $T\colon k \mapsto k + 1
\pmod{8}$.
\end{definition}

\begin{remark}
The number 8 is not arbitrary.  The dimension theorem (T3 in
the RS derivation chain) forces $D = 3$, and the binary nature
of ledger entries forces $2^D = 8$ phases per epoch.
\end{remark}

%% ============================================================
\section{Why $\mathbb{R}$ is Insufficient}
\label{sec:real}
%% ============================================================

\begin{theorem}[$\mathbb{R}$ Cannot Represent $C_8$ Faithfully]
\label{thm:real-fail}
There is no faithful group homomorphism
$\rho\colon C_8 \to GL_1(\mathbb{R}) = \mathbb{R}^\times$.
\end{theorem}

\begin{proof}
If $\rho\colon C_8 \to \mathbb{R}^\times$ is a group
homomorphism and $g$ is the generator of $C_8$ with
$g^8 = 1$, then $\rho(g)^8 = 1$.  The only real solutions
to $x^8 = 1$ are $x = \pm 1$.  If $\rho(g) = 1$, then
$\rho$ is trivial.  If $\rho(g) = -1$, then the kernel
contains $g^2$, so $\rho$ is not faithful.  In either case,
the 8 distinct phases cannot be distinguished.
\end{proof}

\begin{corollary}[$\mathbb{C}$ is the Minimum Extension]
\label{cor:cmin}
$\mathbb{C} = \mathbb{R}[x]/(x^2 + 1)$ is the smallest
field extension of $\mathbb{R}$ over which $C_8$ has a
faithful 1-dimensional representation.
\end{corollary}

\begin{proof}
Over $\mathbb{C}$, the map $g \mapsto \omega = e^{i\pi/4}$
is a faithful representation, since $\omega^k \neq 1$ for
$k = 1, \ldots, 7$.  No smaller extension of $\mathbb{R}$
suffices, because any algebraic extension $\mathbb{R}[\alpha]$
where $\alpha$ is an 8th root of unity generates at least
$\mathbb{R}[i]$ (since $\omega^2 = i$).
\end{proof}

%% ============================================================
\section{The Imaginary Unit}
\label{sec:i}
%% ============================================================

\begin{definition}[Tick-Advance Operator]
\label{def:omega}
The tick-advance operator is the primitive 8th root of unity:
\begin{equation}
  \omega = e^{i\pi/4}
  = \frac{1+i}{\sqrt{2}},
  \qquad \omega^8 = 1.
\end{equation}
Each tick advances the phase by $\pi/4$ ($45^\circ$).
\end{definition}

\begin{theorem}[Imaginary Unit from 8-Tick Epoch]
\label{thm:i}
The quarter-period operator $i = \omega^2$ satisfies
$i^2 = -1$.
\end{theorem}

\begin{proof}
A quarter-period is $8/4 = 2$ ticks, so the quarter-period
operator is $\omega^2 = e^{i\pi/2} = i$.  Squaring:
\begin{equation}
  i^2 = \omega^4 = e^{i\pi} = -1.
\end{equation}
Two quarter-periods (4 ticks) produce the half-period,
which is the phase reversal operator $-1$.
\end{proof}

\begin{remark}
This derivation does not assume complex numbers \emph{a priori}.
Starting from the cyclic group $C_8$ and asking for the
smallest field over which all its irreducible representations
exist, one is forced to adjoin a root of $x^2 + 1 = 0$ to
$\mathbb{R}$.  Calling this root $i$ recovers the standard
complex number field.
\end{remark}

%% ============================================================
\section{The 8 Roots of Unity}
\label{sec:roots}
%% ============================================================

\begin{proposition}[8-Tick Roots]
\label{prop:roots}
The 8 tick phases, in their complex representation, are:
\begin{equation}
  \omega_k = e^{ik\pi/4}, \quad k = 0, 1, \ldots, 7.
\end{equation}
Explicitly:
\begin{center}
\renewcommand{\arraystretch}{1.3}
\begin{tabular}{cccl}
\toprule
$k$ & $\omega_k$ & Phase & Role \\
\midrule
0 & $1$ & $0^\circ$ & Ground state \\
1 & $(1+i)/\sqrt{2}$ & $45^\circ$ & Tick advance \\
2 & $i$ & $90^\circ$ & Quarter period \\
3 & $(-1+i)/\sqrt{2}$ & $135^\circ$ & \\
4 & $-1$ & $180^\circ$ & Half period (reversal) \\
5 & $(-1-i)/\sqrt{2}$ & $225^\circ$ & \\
6 & $-i$ & $270^\circ$ & Three-quarter period \\
7 & $(1-i)/\sqrt{2}$ & $315^\circ$ & \\
\bottomrule
\end{tabular}
\end{center}
\end{proposition}

\begin{remark}[Connection to $\alpha^{-1}$]
The gap weight $w_8$ that determines the fine-structure
constant $\alpha^{-1} \approx 137.036$ is constructed as a
sum over all 8 roots of unity, weighted by the $\varphi$-ladder
coupling.  The fine-structure constant thus encodes the
$\mathbb{C}$-structure of the 8-tick epoch.
\end{remark}

%% ============================================================
\section{Euler's Identity}
\label{sec:euler}
%% ============================================================

\begin{theorem}[Euler's Identity as Cycle Completion]
\label{thm:euler}
Euler's identity $e^{i\pi} + 1 = 0$ is the 8-tick
cycle-completion identity: the half-period amplitude plus the
ground state equals zero.
\end{theorem}

\begin{proof}
The ground state (0 ticks) is $\omega^0 = 1$.  The half-period
(4 ticks) is $\omega^4 = -1$.  Their sum:
\begin{equation}
  \omega^4 + \omega^0 = -1 + 1 = 0.
\end{equation}
Since $\omega^4 = e^{i \cdot 4\pi/4} = e^{i\pi}$, this is
$e^{i\pi} + 1 = 0$.
\end{proof}

\begin{remark}
Euler's identity is sometimes called the ``most beautiful
equation in mathematics'' because it connects the five
fundamental constants $e$, $i$, $\pi$, $0$, and $1$.  In RS,
this beauty has a structural explanation: the identity is the
statement that opposite phases in the 8-tick cycle cancel
perfectly.  This is a consequence of the $C_8$ group structure,
not a coincidence of analysis.
\end{remark}

\begin{corollary}[Phase Cancellation]
\label{cor:cancel}
More generally, opposite phases in the 8-tick cycle always
cancel:
\begin{equation}
  \omega^k + \omega^{k+4} = 0 \quad
  \text{for all } k \in \mathbb{Z}.
\end{equation}
This is the RS origin of destructive interference.
\end{corollary}

%% ============================================================
\section{The 2D Representation}
\label{sec:2d}
%% ============================================================

\begin{proposition}[Matrix Representation]
\label{prop:matrix}
Over $\mathbb{R}$, the faithful irreducible representation
of the tick-advance operator is the $2 \times 2$ rotation
matrix
\begin{equation}
  R(\pi/4) = \begin{pmatrix}
    \cos(\pi/4) & -\sin(\pi/4) \\
    \sin(\pi/4) & \cos(\pi/4)
  \end{pmatrix} = \frac{1}{\sqrt{2}}
  \begin{pmatrix} 1 & -1 \\ 1 & 1 \end{pmatrix}.
\end{equation}
The quarter-period operator $i$ corresponds to
\begin{equation}
  R(\pi/2) = \begin{pmatrix} 0 & -1 \\ 1 & 0 \end{pmatrix},
  \qquad R(\pi/2)^2 = -I,
\end{equation}
confirming $i^2 = -1$ as a matrix identity.
\end{proposition}

\begin{remark}
This is why quantum mechanics requires a 2D internal
structure (complex amplitudes): the state space must
accommodate the 8-tick phase rotation, which is irreducibly
2-dimensional over $\mathbb{R}$.
\end{remark}

%% ============================================================
\section{Mathematical Constants as RS-Native}
\label{sec:native}
%% ============================================================

\begin{center}
\renewcommand{\arraystretch}{1.3}
\begin{tabular}{lll}
\toprule
\textbf{Constant} & \textbf{RS Origin}
  & \textbf{Derivation} \\
\midrule
$\varphi = (1+\sqrt{5})/2$ & Self-similar ledger (T6) &
  Unique fixed point of $J$-cost \\
$i$ \;($i^2 = -1$) & 8-tick quarter period &
  $\omega^2$ where $\omega^8 = 1$ \\
$\pi$ & Half-period angle &
  $\omega^4 = e^{i\pi}$ \\
$e$ & Recognition partition function &
  $\sum_n 1/n!$ from path counting \\
$e^{i\pi} + 1 = 0$ & Cycle completion &
  $\omega^4 + \omega^0 = 0$ \\
$\ln\varphi = k_R$ & J-bit &
  RS Boltzmann constant \\
\bottomrule
\end{tabular}
\end{center}

\begin{remark}
The key claim is not that RS ``explains'' these constants
\emph{ex nihilo}---they are theorems of mathematics regardless
of physics.  The claim is that the same structure
($J$-cost $\to$ $\varphi$ $\to$ $D = 3$ $\to$ 8-tick
$\to$ $\mathbb{C}$) that forces the physical laws also forces
the mathematical structures.  Mathematics is not imported into
physics; it is an inevitable consequence of the same
recognition structure.
\end{remark}

%% ============================================================
\section{Predictions and Falsifiers}
\label{sec:predictions}
%% ============================================================

\begin{enumerate}
  \item \textbf{Complex quantum mechanics is necessary.}
  Quantum mechanics cannot be formulated with real Hilbert
  spaces alone.  The recent experimental results of
  Renou et al.~\cite{Renou2021} confirm this, consistent
  with the RS prediction.

  \item \textbf{No quaternionic quantum mechanics.}
  The 8-tick structure requires $\mathbb{C}$, not
  $\mathbb{H}$ (quaternions).  Quaternionic quantum
  mechanics would require a 16-tick or higher epoch,
  inconsistent with $2^D = 8$.

  \item \textbf{Falsifier.}
  Discovery of a physical process requiring a field
  extension of $\mathbb{C}$ (e.g., quaternionic phases)
  would falsify the 8-tick epoch structure.
\end{enumerate}

%% ============================================================
\section{Conclusion}
\label{sec:conclusion}
%% ============================================================

Complex numbers are forced by the 8-tick epoch of Recognition
Science.  The cyclic group $C_8$ generated by the ledger's
update cycle cannot be faithfully represented over
$\mathbb{R}$; the minimum extension is $\mathbb{C}$.  The
imaginary unit $i = \omega^2$ is the quarter-period operator,
$i^2 = -1$ is the half-period reversal, and Euler's identity
$e^{i\pi} + 1 = 0$ is the cycle-completion identity.
Mathematics is not imported into RS physics---it is the same
structure, viewed from the formal side.

%% ============================================================
\begin{thebibliography}{99}

\bibitem{Renou2021}
M.-O.~Renou \textit{et al.}, ``Quantum theory based on real
numbers can be experimentally falsified,'' Nature
\textbf{600}, 625--629 (2021).

\bibitem{Procopio2022}
L.~M.~Procopio \textit{et al.}, ``Experimental test of local
observer-independence,'' Sci.\ Adv.\ \textbf{5}, eaaw9832
(2019).

\bibitem{Penrose2004}
R.~Penrose, \emph{The Road to Reality}, Jonathan Cape (2004).

\end{thebibliography}

\end{document}
